\documentclass{beamer}

\usepackage[italian]{babel}
\usepackage[utf8]{inputenc}
\usepackage{listings}
\usepackage{xcolor}

\usetheme{Madrid}

\title{Analisi della Struttura di una Web App PCTO}
\author{Prof. Fedeli Massimo}
\date{}

\lstset{
	language=HTML,
	basicstyle=\ttfamily\small,
	keywordstyle=\color{blue},
	stringstyle=\color{red},
	breaklines=true
}

\begin{document}
	
	\frame{\titlepage}
	
	%------------------------------------------------
	
	\begin{frame}
		\frametitle{Obiettivo del progetto}
		
		Il progetto realizza una semplice applicazione web per la gestione
		delle esperienze di PCTO degli studenti.
		
		Le \textbf{funzionalità} principali sono:
		
		\begin{itemize}
			\item autenticazione degli utenti
			\item inserimento di una esperienza PCTO
			\item ricerca delle esperienze per intervallo di date
			\item salvataggio dei dati in un database MySQL
		\end{itemize}
		
		L'applicazione è composta da due parti principali:
		
		\begin{itemize}
			\item Frontend (HTML + JavaScript)
			\item Backend (PHP + MySQL)
		\end{itemize}
		
	\end{frame}
	
	%------------------------------------------------
	
	\begin{frame}
		\frametitle{Architettura dell'applicazione}
		
		L'applicazione segue una architettura client-server.
		
		\bigskip
		
		\textbf{Client (browser)}
		
		\begin{itemize}
			\item interfaccia HTML
			\item codice JavaScript
			\item invio delle richieste HTTP
		\end{itemize}
		
		\bigskip
		
		\textbf{Server}
		
		\begin{itemize}
			\item API scritte in PHP
			\item gestione delle sessioni
			\item accesso al database MySQL
		\end{itemize}
		
	\end{frame}
	
	%------------------------------------------------
	
	\begin{frame}
		\frametitle{Struttura dei file del progetto}
		
		Il progetto è organizzato in più file con responsabilità diverse.
		
		\bigskip
		
		\textbf{Frontend}
		
		\begin{itemize}
			\item \textbf{index.html}
		\end{itemize}
		
		\bigskip
		
		\textbf{Backend}
		
		\begin{itemize}
			\item config.php
			\item db.php
			\item auth.php
		\end{itemize}
		
		\bigskip
		
		\textbf{API}
		
		\begin{itemize}
			\item api\_login.php
			\item api\_logout.php
			\item api\_session.php
			\item api\_experiences.php
		\end{itemize}
		
	\end{frame}
	
	%------------------------------------------------
	
	\begin{frame}[fragile]
		\frametitle{Il Frontend: index.html}
		
		Il file \texttt{index.html} contiene l'interfaccia dell'applicazione.
		
		Gli elementi principali sono:
		
		\begin{itemize}
			\item form di login
			\item form per inserire una esperienza
			\item ricerca delle esperienze
		\end{itemize}
		
		\begin{lstlisting}
			<input id="username" placeholder="username">
			<input id="password" type="password">
			<button onclick="login()">Login</button>
		\end{lstlisting}
		
		Quando l'utente preme il pulsante viene chiamata
		una funzione JavaScript.
		
	\end{frame}
	
	%------------------------------------------------
	
	\begin{frame}[fragile]
		\frametitle{Comunicazione con il server}
		
		Il frontend comunica con il backend tramite
		richieste HTTP usando la funzione \textbf{fetch}().
		
		\begin{lstlisting}
			function api(url,method,data){
				
				let options = {
					method:method,
					headers:{'Content-Type':'application/json'},
					credentials:'same-origin'
				}
				
				if(data){
					options.body = JSON.stringify(data)
				}
				
				return fetch(url,options)
			}
		\end{lstlisting}
		
		Questa funzione invia richieste alle API del server.
		
	\end{frame}
	
	%------------------------------------------------
	
	\begin{frame}
		\frametitle{Backend in PHP}
		
		Il backend è composto da diversi file PHP che svolgono
		funzioni specifiche.
		
		\bigskip
		
		\textbf{config.php}
		
		\begin{itemize}
			\item contiene le configurazioni
			\item definisce i parametri del database
			\item avvia la sessione
		\end{itemize}
		
		\bigskip
		
		\textbf{db.php}
		
		\begin{itemize}
			\item gestisce la connessione al database
			\item utilizza PDO
		\end{itemize}
		
	\end{frame}
	
	%------------------------------------------------
	
	\begin{frame}[fragile]
		\frametitle{Connessione al database}
		
		Il file \textbf{db.php }crea una connessione PDO
		al database MySQL.
		
		\begin{lstlisting}
			$dsn = 'mysql:host=' . DB_HOST .
			';dbname=' . DB_NAME .
			';charset=utf8mb4';
			
			$pdo = new PDO($dsn, DB_USER, DB_PASS, $options);
		\end{lstlisting}
		
		PDO è una libreria che permette di lavorare con i database
		in modo sicuro e portabile.
		
	\end{frame}
	
	%------------------------------------------------
	
	\begin{frame}
		\frametitle{Gestione dell'autenticazione}
		
		Il file \textbf{auth.php }contiene funzioni comuni
		usate dalle API.
		
		Tra le principali:
		
		\begin{itemize}
			\item invio delle risposte JSON
			\item verifica che l'utente sia autenticato
			\item lettura dei dati JSON ricevuti
		\end{itemize}
		
		Questo permette di evitare duplicazione del codice.
		
	\end{frame}
	
	%------------------------------------------------
	
	\begin{frame}
		\frametitle{API REST}
		
		Le API sono programmi PHP che ricevono richieste HTTP.
		
		Ogni API svolge una funzione specifica.
		
		\bigskip
		
		\textbf{api\_login.php}
		
		gestisce l'autenticazione dell'utente.
		
		\bigskip
		
		\textbf{api\_logout.php}
		
		termina la sessione.
		
		\bigskip
		
		\textbf{api\_experiences.php}
		
		gestisce inserimento e ricerca delle esperienze.
		
	\end{frame}
	
	%------------------------------------------------
	
	\begin{frame}[fragile]
		\frametitle{Login dell'utente}
		
		Durante il login il server controlla username e password.
		
		\begin{lstlisting}
			$stmt = $pdo->prepare(
			'SELECT id, username, password_hash
			FROM utenti WHERE username = ?'
			);
			
			$stmt->execute([$username]);
		\end{lstlisting}
		
		La password viene verificata con
		
		\begin{lstlisting}
			password_verify($password,
			$user['password_hash']);
		\end{lstlisting}
		
		Questo garantisce maggiore sicurezza.
		
	\end{frame}
	
	%------------------------------------------------
	
	\begin{frame}
		\frametitle{Gestione delle sessioni}
		
		Dopo il login il server salva i dati dell'utente
		nella sessione.
		
		\bigskip
		
		Questo permette di:
		
		\begin{itemize}
			\item riconoscere l'utente nelle richieste successive
			\item proteggere le API
		\end{itemize}
		
		Le API possono verificare se l'utente è autenticato
		prima di eseguire le operazioni.
		
	\end{frame}
	
	%------------------------------------------------
	
	\begin{frame}[fragile]
		\frametitle{Inserimento di una esperienza}
		
		L'inserimento avviene tramite una richiesta POST.
		
		\begin{lstlisting}
			INSERT INTO esperienze_pcto
			(id_studente, ente_ospitante,
			mansione, data_inizio, data_fine,
			ore_totali)
			VALUES (?, ?, ?, ?, ?, ?)
		\end{lstlisting}
		
		I parametri sono passati tramite JSON
		dal frontend.
		
	\end{frame}
	
	%------------------------------------------------
	
	\begin{frame}
		\frametitle{Ricerca delle esperienze}
		
		L'API permette anche di recuperare
		le esperienze per intervallo di date.
		
		La query SQL utilizza una JOIN tra
		le tabelle:
		
		\begin{itemize}
			\item esperienze\_pcto
			\item studenti
		\end{itemize}
		
		In questo modo vengono restituiti
		anche nome e cognome dello studente.
		
	\end{frame}
	
	%------------------------------------------------
	
	\begin{frame}
		\frametitle{Flusso completo dell'applicazione}
		
		1. L'utente apre la pagina web
		
		2. Inserisce username e password
		
		3. Il browser invia una richiesta all'API di login
		
		4. Il server verifica le credenziali
		
		5. Se il login è valido viene creata una sessione
		
		6. L'utente può inserire o consultare esperienze PCTO
		
	\end{frame}
	
	%------------------------------------------------
	
	\begin{frame}
		\frametitle{Conclusioni}
		
		Questo progetto mostra come realizzare
		una semplice applicazione web completa.
		
		Sono stati utilizzati:
		
		\begin{itemize}
			\item HTML
			\item JavaScript
			\item PHP
			\item MySQL
		\end{itemize}
		
		Il progetto introduce concetti fondamentali
		dello sviluppo web:
		
		\begin{itemize}
			\item architettura client-server
			\item API REST
			\item autenticazione
			\item accesso al database
		\end{itemize}
		
	\end{frame}
	
\end{document}